\section{Necessary product order for full-vector common partitions}\label{sec:common-rank}
We prove Theorem~\ref{thm:common-rank} by an explicit positive diagonal
family. The proof concerns all vector fields, as in Theorem~\ref{thm:matrix},
rather than the smaller coexact domain of the endpoint lower theorem.
It uses only local fields and finite-dimensional rank-nullity.

\subsection{The coefficient family and every local channel}
Let $n=3m+3$ and choose centers
\[
 x_j=(3j+1,1,\ldots,1),\qquad 0\le j<m.
\]
They lie away from the periodic wrap boundary. With
$\sigma=\alpha_D^{1/P}c$, define
\[
 a_0=\min\{c/(8D),\sigma/2\}>0,\qquad
 M=r n^D,\qquad H=\sigma M>\sigma.
\]
In the canonical direction/component basis $(i,k)$ of $\R^{Dh}$ let
$A(x)$ be diagonal, with entries
\begin{equation}\label{eq:rank-profile}
 \lambda_{i,k}(x)=
 \begin{cases}
 a_0,&k=1\text{ and }x\in\{x_j,x_j-e_i\}\text{ for some }j,\\
 H,&\text{otherwise}.
 \end{cases}
\end{equation}
At a center all $D$ direction channels for component one are weak; at an
incoming start only its specified direction/component channel is weak.
All eigenvalues are positive. For each fixed direction the starts $x_j$
and $x_j-e_i$ are distinct. Across centers their first coordinates lie in
disjoint sets $\{3j,3j+1\}$. Channels with different direction labels remain
distinct even when they share a start. There are therefore exactly $2Dm$
weak channels.

Since $a_0\le\sigma/2$, each weak channel has capped cost one in
\eqref{eq:cap}. Every other channel has capped cost
$\alpha_Dc^PH^{-P}=M^{-P}$. Thus, including all strong channels,
\begin{equation}\label{eq:rank-budget}
 F_c(A)=2Dm+(M-2Dm)M^{-P},\qquad
 2Dm\le F_c(A)\le2Dm+1.
\end{equation}
Indeed $0\le M-2Dm\le M$ and $P>1$, so the tail is at most
$M^{1-P}\le1$. Increasing $h$ increases $H$ and keeps this strong tail
bounded; it introduces no uncontrolled $h$-sized budget contribution.
The choice of $H$ is permitted by the arbitrary finite profile hypothesis.

\subsection{The actual low space and the restriction rank}
Let $U_j(x)=1_{\{x_j\}}(x)e_1$, where $e_1\in\R^h$ is the first
canonical unit vector. These fields are orthonormal in counting units.
Their only nonzero direction-labelled forward differences are $-e_1$ at
$x_j$ and $+e_1$ at $x_j-e_i$, for each direction $i$. Both original
starts are assigned coefficient $a_0$ by \eqref{eq:rank-profile}; there
are no further wrap contributions. Consequently
\[
 Q_A(U_j)=2Da_0\le c/4.
\]
No nearest-neighbor edge joins distinct centers, and their gradient-start
supports are disjoint. Hence the quadratic cross terms vanish and
\begin{equation}\label{eq:rank-lowspace}
 L_m=\operatorname{span}_{\R}\{U_0,\ldots,U_{m-1}\},\quad
 \dim L_m=m,\quad
 Q_A(U)\le(c/4)\sum_x\|U(x)\|^2\quad(U\in L_m).
\end{equation}
These are general vector functions. No divergence or coexact property is
asserted for them.

For any common partition $\mathcal R$ into $R$ nonempty rectangles, its
actual rows on this space are
\[
 q_B\left(\sum_jt_jU_j\right)
 =V^{-1/2}|B|^{-1/2}
       \left(\sum_{j:x_j\in B}t_j\right)e_1.
\]
Each rectangle supplies at most one scalar functional on $L_m$, so
$\rank(C_{\mathcal R}|_{L_m})\le R$. If $R<m$, rank-nullity gives a
nonzero $U\in L_m$ with $C_{\mathcal R}U=0$. Coercivity with any finite
quadratic penalty, after multiplication by $V$, would imply
\[
 c\sum_x\|U(x)\|^2\le Q_A(U)
                    \le(c/4)\sum_x\|U(x)\|^2,
\]
a contradiction. Thus every successful common partition has $R\ge m$.
The same argument works for any penalty on these linear observations with
$\Phi(0)=0$, without a differentiability assumption on that penalty.

The full field-space rank is exactly $hR$, not merely a padded output
dimension. To see this, for each $B$ choose a site $z\in B$. A field
supported at $z$ with value $\sqrt{V|B|}\,e_k$ has observation $e_k$ in
row $B$ and zero in every other row. These fields show surjectivity onto
all $hR$ row coordinates. Therefore necessary genuine rank is at least
$hm$. Since $F_c(A)\le2Dm+1\le(2D+1)m$, this proves
\eqref{eq:product-lower}.

\subsection{Joint growth and the measurement-class restriction}
Set $h=m$. The necessary rank is at least $m^2$, while
$h+F_c(A)\le(2D+1)m+1$. Their ratio tends to infinity for fixed $D,c$.
Thus a dimension-only $C_D(h+F_c(A))$ bound fails for these common
full-vector partitions, even with arbitrary finite penalty. The product
lower bound and Theorem~\ref{thm:matrix} have matching $hF_c(A)$ order
along this family, since $F_c(A)\ge2Dm$ makes $1+F_c(A)$ comparable to
$F_c(A)$. Sharp coefficients, additive constants and per-profile optimality
are not determined.

The extra factor $h$ measures the cost of retaining every component mean
on every common partition member; it does not count $h$ independent tests
of the one-component low space. The argument does not apply to distinct
componentwise partitions, projected or channel-adapted rows, or compressed
postprocessed outputs. It requires no external theorem about the rank of
those broader observation classes. These elementary rank and local-field
arguments carry no claim of historical priority, a new general principle,
formal verification, or physical transfer.
